\subsection{Excluding order $91$}\label{sec:decic91}
We prove that $X_1(91)$ has no noncuspidal effective rational divisor of degree $10$.
Put $X=X_1(91)$. The exact winding calculation shows that the character of
every winding-zero constituent is trivial on $3$. The winding-quotient
argument therefore gives
\[
 (\langle3\rangle-1)J_1(91)(\mathbb Q)\quad\text{finite}.
\]
We do not assume that $J_1(91)(\mathbb Q)$ has rank zero.

We use the normalization
\[
 Y=X_1(7)\mathbin{\times}_{X(1)}X_0(13),\qquad
 Z=X_0(91).
\]
The map $X\longrightarrow Y$ sends $(E,P)$ to
$(E,13P,\langle7P\rangle)$ and has degree $12$.
The curve $Y$ has genus $23$. Its diamond group over $Z$ has order $3$;
write $\tau$ for a generator and $\pi:Y\longrightarrow Z$ for the quotient.
Since $\langle3\rangle$ induces a generator on $Y$, the group
$(\tau-1)J_Y(\mathbb Q)$ is finite.
All these curves have good reduction at $2$.

\begin{lemma}\label{lem:decic91-relative-formal-fibre}
Let $R=\mathbb Z_2$, and let $Y/R$ be a smooth proper modular curve of odd level.
Let $\tau$ be a diamond automorphism of order $3$, put $Z=Y/\langle\tau\rangle$,
and write $\pi:Y\longrightarrow Z$ for the quotient map.
Assume that $Z/R$ is smooth.
Suppose $(\tau-1)J_Y(\mathbb Q)$ is finite.
Let $A$ be an effective rational divisor of degree $d$ on the generic fibre.
Let $C$ be an effective divisor over $R$ with $\overline A=\overline C$, and suppose
$(\tau-1)[C-dc_0]$ is torsion in $J_Y(\mathbb Q_2)$, where $c_0$ is a rational cusp.
Write
\[
 C=C_0+\pi^*E,\qquad \deg E=k,
\]
and suppose the map $E'\mapsto C_0+\pi^*E'$ is a closed formal immersion
at $\overline E$. If the restriction of
\[
 U=(\operatorname{Cartier}-1)(\tau^*-1)H^0(Y_{\mathbb F_2},\Omega^1)
\]
to $\overline C$ has rank $d-k$, then $A$ contains $C_0$ as a divisor over $\mathbb Q_2$.
\end{lemma}
\begin{proof}
Put $\alpha=\tau-1$ and $M=T_2-\langle2\rangle$ on the full Jacobian $J_Y$.
The class $\alpha[A-C]$ is torsion and has zero reduction.
Lemma~\ref{lem:dyadic-formal-Hecke-kernel} gives $M\alpha[A-C]=0$.
Thus $A$ belongs to the formal fibre through $C$ of
\[
 \Psi:Y^{(d)}\longrightarrow J_Y,\qquad
 D\longmapsto M\alpha[D-dc_0].
\]
By Eichler--Shimura, the cotangent image of $\Psi$ is $U$;
the invertible diamond factor does not change this image.

Choose $d-k$ local coordinates of $J_Y$ whose pullbacks have independent
differentials at $\overline C$.
Their common fibre is formally smooth over $R$ of relative dimension $k$.
The formal family $C_0+\pi^*Z^{(k)}$ lies in the full fibre of $\Psi$,
since $(\tau-1)\pi^*=0$.
Its tangent map to the chosen smooth fibre is an isomorphism between
$k$-dimensional spaces. It is therefore a formal isomorphism.
The full fibre equals this family, proving the assertion.
\end{proof}

We use the lemma only in the following two situations, where the stated
closed immersion can be checked directly. At an unramified orbit of three
cusps, subtract the minimum of their multiplicities. At least one residual
multiplicity is zero, and the divisor at that point recovers the monic
polynomial defining the moving divisor downstairs. The same argument applies
to a cubic cusp whose three geometric points form one orbit.
At a tame fixed quadratic point $P$, choose, after unramified base change,
a parameter $z$ with $\tau z=\zeta_3z$. The quotient parameter is $z^3$.
Near $P+\pi^*(\pi(P))=4P$, the family is given by
$z(z^3-a)$, whose coefficient of $z$ recovers $a$.
These descriptions give closed formal immersions over the split unramified
extension, and hence over $R$ by descent.


We give an explicit model for the computations. Over $\mathbb F_2(t)$ put
\[
 b=t^3-t^2,\quad c=t^2-t,\quad
 E_t:y^2+(1-c)xy-by=x^3-bx^2.
\]
The point $(0,0)$ has order $7$. Put $j_7(t)=j(E_t)$ and
\[
 j_{13}(s)=
 \frac{(s^2+5s+13)(s^4+7s^3+20s^2+19s+1)^3}{s}.
\]
The normalization of $j_{13}(s)=j_7(t)$ is $Y_{\mathbb F_2}$.
Its function field has degree $14$ over $\mathbb F_2(t)$ and genus $23$.
A generator acts by $\tau(t)=1/(t+1)$ and $\tau(s)=s$.
The matrices of Cartier and $\tau^*$ have ranks
\[
 \operatorname{rank}(\operatorname{Cartier})=21,\qquad
 \operatorname{rank}(\tau^*-1)=16,\qquad \dim U=14.
\]
These statements are checked in the native function field, together with
$\tau^3=1$ and commutation with Cartier.

\begin{proposition}\label{prop:decic91-exclusion}
An elliptic curve over a number field of degree $10$ has no point of order $91$.
\end{proposition}
\begin{proof}
Suppose $D$ is a noncuspidal effective rational divisor of degree $10$ on $X$.
The complete census at $2$ gives no noncuspidal closed point of degree at most $7$,
three of degree $8$, $144$ of degree $9$, and none of degree $10$.
Thus the noncuspidal part of $\overline D$ is empty, one degree-$8$ point,
or one degree-$9$ point.
Write $A$ for the pushforward of $D$ to $Y$; it is entirely noncuspidal and has degree $10$.

Suppose first that $\overline D$ is entirely cuspidal.
The source cusp signatures over $\mathbb Q$ and $\mathbb F_2$ are respectively
\[
 1^{36}6^6 12^3 36,
 \qquad 1^{36}3^{12}12^6.
\]
Their possible degree-$10$ pushforwards to $Y$ are supported on three rational
cusps $c_1,c_2,c_3$ of width $91$ and one cubic cusp $q$ of width $13$.
All four closed cusps are globally intact at $2$.
The automorphism $\tau$ cycles the three rational cusps and acts transitively
on the geometric points of $q$.
The primitive-column calculation verifies every source cusp norm and the
complete list of $120$ vectors
\[
 C=a_1c_1+a_2c_2+a_3c_3+bq,
 \qquad a_1+a_2+a_3+3b=10.
\]
Put $m=\min(a_1,a_2,a_3)$ and $k=m+b$.
The maximal pullback part of $C$ is
\[
 \pi^*E=m(c_1+c_2+c_3)+bq,
 \qquad \deg E=k.
\]
Hence $C_0=C-\pi^*E$ has degree $10-3k>0$ and is cuspidal.
For every one of the $120$ vectors, the exact restriction rank of $U$ is
\[
 \operatorname{rank}(U|_C)=10-k.
\]
The quotient is unramified at these cusps. Lemma~\ref{lem:decic91-relative-formal-fibre}
therefore implies that $A$ contains $C_0$, a contradiction.

Next suppose $\overline D=P_8+R_2$.
The three possible $P_8$ have $j=0$ and come from two twists over $\mathbb F_{256}$,
each with cyclic group of order $273$.
The order-$3$ elliptic automorphism sends $P$ to $16P$ for the retained
order-$91$ generator. It therefore sends $13P$ to $2(13P)$ and preserves
$\langle7P\rangle$. Thus the image on $Y$ is geometrically fixed by $\tau$.
The fixed locus on $Y_{\mathbb F_2}$ consists of two quadratic points;
consequently the image of $P_8$ is $4p$, for one of these points $p$.
The complement $R_2$ pushes forward to an effective degree-$2$ divisor
on $\{c_1,c_2,c_3\}$.

At each $p$, the tangent eigenvalue of $\tau$ has order $3$.
The fixed locus therefore gives a unique unramified quadratic lift $P$
over $\mathbb Z_2$, with $\tau P=P$.
Put $C=4P+R_2$. The class $(\tau-1)[C-10c_0]$ is cuspidal torsion.
The maximal pullback part is $3P=\pi^*\pi(P)$; its degree downstairs is $2$.
For both choices of $p$ and all six possibilities for $R_2$, the exact rank is
\[
 \operatorname{rank}(U|_{4p+\overline R_2})=8=10-2.
\]
The tame fixed-point case of Lemma~\ref{lem:decic91-relative-formal-fibre}
implies that $A$ contains $P+R_2$. In particular it contains a cusp, a contradiction.

It remains to exclude $\overline D=P_9+\overline c$, where $c$ is a rational cusp on $X$.
Put $\sigma=\langle3\rangle^{-1}$ and $A_3=T_3-1-3\langle3\rangle$.
The operator $A_3$ kills rational torsion. Moreover, for every rational cusp $c$,
\[
 T_3c=c+3\langle3\rangle c,
\]
so $A_3c=0$ as a divisor.
Applying $A_3(\sigma-1)$ to $[D-10c_0]$ gives the necessary principal divisor
\[
 B=\sigma T_3P_9-T_3P_9-\sigma P_9-2P_9+3\langle3\rangle P_9.
\]
This test is independent of the total degree of $D$ once the rational cusp
complement has been annihilated.

The corrected canonical-ring certificate for $X_1(91)$ checks these same
$144$ degree-$9$ closed points. It retains four source representatives,
all $36$ simultaneous target diamond translates and both point conventions,
giving $288$ tests that include every necessary $B$.
The canonical space has dimension $265$, and its ordered products span
$H^0(X,2K_X)$ of dimension $792$.
For each $B=B^+-B^-$, nonzero canonical sections $F,H$ vanish on $B^+$.
If $f\in\mathcal L(B)$, then $G=fF$ vanishes on $B^-$ and satisfies
$GH\in FH^0(X,K_X)$.
The resulting map
\[
 H^0(X,K_X)\longrightarrow
 H^0(B^-,K_X|_{B^-})\oplus
 H^0(X,2K_X)/FH^0(X,K_X)
\]
is injective in every test: all retained $265\times590$ matrices have rank $265$.
The source and target points, marked isogenies, canonical identities and jets
are part of the exact certificate, with independent finite-field and binary
matrix checks. Hence $B$ is not principal. This excludes the last case.
\end{proof}

The degree-$10$ evidence is stored in
\path{continuation_20261030c/C91}.
The native packet contains the defining field, exact two-generator prime ideals,
Cartier and diamond matrices, $33$ individual cusp restriction kernels,
and $12$ mixed CM restriction kernels.
Independent binary elimination checks all $120$ relative cusp ranks,
all $12$ CM ranks and all $288$ inherited canonical ranks.
A second native replay reconstructs the field, all six prime ideals and all
$45$ differential kernels. It also checks
$\Phi_{13}(j_{13}(s),j_{13}(1/s))=0$, recomputes Cartier and the diamond action,
and verifies all relative ranks.
The inherited canonical-ring data are copied with their source hashes;
only their degree-$9$ nonprincipality statement is used here.
